\documentclass[11pt]{article}

% Prefix for numedquestion's
\newcommand{\questiontype}{Question}


% Use this if your "written" questions are all under one section
% For example, if the homework handout has Section 5: Written Questions
% and all questions are 5.1, 5.2, 5.3, etc. set this to 5
% Use for 0 no prefix. Redefine as needed per-question.
\newcommand{\writtensection}{0}

\usepackage{amsmath, amsfonts, amsthm, amssymb}  % Some math symbols
\usepackage{mathtools}

\usepackage{centernot}
\usepackage{mathtools}

\setlength{\parindent}{0pt}

\begin{document}

\section{One-Way Permutations Imply UOWHF}
From the end of lecture 10 we saw that the following hash function is $2-$wise independent.
$$h_{a,b}(x) = (ax+b)_{1,n}$$
We extend this idea to create a UOWHF based on one way permutation $f$. First, we define the function $chop(x)$ to cut down a length $n$ string to length $n-1$ by deleting the first bit. The hash function $h$ is defined as follows:
$$h_{a,b}(x) = chop(a(f(x))+b)$$
\textbf{Proof of Security:} To demonstrate that this protocol is secure, we need to argue that an adversary that could break the Universal One-Way property of this function could break the one way permutation property of $f$. Consider an adversary $A$ that could find another $x'$ such that:
$$h_{a,b}(x) = chop(a(f(x'))+b)$$
Observe that there is only one other $x'$ value that could map to the same $h_{a,b}(x)$ value. An adversary $A'$ could have an advantage over the permutation $f$ as follows.
\begin{enumerate}
    \item Adversary $A$ has the goal of finding inverse of value $y=f(x)$. Given oracle access to adversary $A'$, it receives value $x'$ and solves for $a,b$:
    $$ay + b = a(f(x')) + b$$
    Since $x$ was chosen at random, $y$ must be random as well. Therefore, the values of $a$ and $b$ must be distributed randomly so when $A$ returns hash function $h_{a,b}$ the function is indistinguishable from completely randomly $h$ from $H$.
    \item Pass the hash function $h_{a,b}$ to $A'$ who finds $x$ with non negligible probability. This is because there are only two solutions to $h_{a,b}(x)$ so the only collision it can find for $h$ would be $x'$ exactly. Meaning it can invert $f$ with non-negligible probability, which is a contradiction.
\end{enumerate} 

\subsection{Decreasing Hash Function}
We claim that we can construct a hash function that maps $\{0,1\}^{poly(n)} \rightarrow \{0,1\}^n$ given a family of permutations $f_{poly(n)}, f_{poly(n)-1},..., f_{n}$. The construction follows from the previous hash function, selecting $poly(n)-n$ random $a,b$ values and at each step appling $x_{i-1} = cut(a_if_i(x_i)+b_i)$. The proof of security follows demonstrates that if this is not a secure hash function, then we can invert some $f_k$ with non negligible probability.
\vspace{1em}

\textbf{Proof:} Firstly observe that an adversary $A$ to be able to invert $h$, it needs to find a collision at some level $i$ for function $f_i$. Because of this, it must be able to find a collision for some function $f_i$ with probability $\epsilon/n$. We can invert this $f_i$ by following the previous proof and replacing the output of level $i$ with our new $a_if(x')+b$. Now, if $A$ outputs input to $h$ that results in a collision, we can run the function and with probability at least $\epsilon/poly(n)$ it will invert $f_k$.

\section{Encryption}
\subsection{Public Key Encryption}
The modern definition of a Public Key Exchange (PKE) protocol is a triple of algorithms defined as follows:
\begin{itemize}
    \item $\mathsf{KeyGen}(1^k,R) \rightarrow pK, sK$
    \item $\mathsf{Enc}(m,pK,R) \rightarrow C$
    \item $\mathsf{Dec}(c',pK,sK) \rightarrow m'$
\end{itemize}
The informal definition is that nobody can get $m$ given $c$ without the secret key.

\subsubsection{Goldwasser-Micali}
Goldwasser and Micali published a paper describing properties that a cryptographic protocol should have. Among these, there is the notion of knowing nothing about the plaintext given the ciphertext. It should be the case that even an encryption of a single bit should be indistinguishable whether its a 0 or 1, which immediately implies that there must be some notion of randomness used.

\subsection{Diffie-Hellman Key Agreement}
\subsubsection{Discrete Log}
The Diffie-Hellman key agreement is based on mathematical properties of prime order groups. Firstly, $\mathbb{Z}_p^*$ is a group of size $p$ of integers $\{1,...,p\}$. If the size of this group is prime, then for any element in this set $g$, we say $g$ generates $\mathbb{Z}_p^*$. This simply means that if we multiply $g$ by itself, we can land on any element in this set. We say for any $a \in \mathbb{Z}_p^*$:
$$a = g^k \mod p$$
for some integer $k$.
\vspace{1em}

The discrete log problem says that for arbitrary $a \in \mathbb{Z}_p^*$, it is hard to find $k$ such that $g^k \mod p = a$.

\subsubsection{Protocol}
The protocol is as follows. 
\begin{enumerate}
    \item Alice announces large prime $p$ and generator $g$ of $\mathbb{Z}_p^*$
    \item Alice and Bob select random elements from the group $a$ and $b$ respsectively and send $A=g^a, B=g^b$ to the other party.
    \item The generated key they share is $g^{ab}$ which can be calculated by taking their selected element and calculating $A^b, B^a$.
\end{enumerate}

\subsection{Semantic Security}
Semantic Security define properties on protocols.
\begin{itemize}
    \item Correctness: Decryption on a valid ciphertexst always produces the message that was encrypted
    \item Security: Given two messages, the probability that any PPT adversary can tell the difference between the produced ciphertexts of the two is negligible.
\end{itemize}
\subsection{El Gamal Public Key Encryption}
The El Gamal Encryption scheme is performed over some large finite group of prime order with chosen generator $g$ as follows:
\begin{enumerate}
    \item $\mathsf{Keygen}(1^n,R) \rightarrow pK=(G,p,g,g^x), sK=x$ where $x$ is a random group element, $G$ is the group, $p$ is the prime order of the group, and $g$ is a generator.
    \item $\mathsf{Enc}(m,pK)$ chooses a random group element $r$ and sends $(g^r, pK^r \cdot m)$
    \item $\mathsf{Dec}(c',sK,pK)$ takes $k = (g^r)^{sK}$ and takes the second part of the ciphertext to generate $m = \frac{pK^r \cdot m}{(g^r)^{sK}} = m$
\end{enumerate}

This scheme is secure under the assumption that decisional Diffie Hellman (DDH) is hard. This problem says that not PPT adversary can tell the following two outputs apart with non-neglible probability:
\begin{enumerate}
    \item $\langle G, p, g, g^x, g^y, g^{xy}\rangle$
    \item $\langle G, p, g, g^x, g^y, g^{r}\rangle$
\end{enumerate}
Where $G$ is a group, $p$ is the size of the group which must be a large prime, $g$ is a generator of the group and $x,y,r$ are random group elements.
\vspace{1em}

\textbf{Proof of Security:} We consider an adversary $A$ that can break the semantic security of this protocol. This means that it must be able to distinguish two ciphertexts with non-negligible probability. The adversary $A'$ that can break $DDH$ given oracle access to $A$ works as follows:
\begin{enumerate}
    \item $A'$ receives $\langle G, p, g, g^x, g^y, g^{c}\rangle$ where $c$ is either $xy$ or a random field element $r$.
    \item $A'$ sends $A$ $(G,p,g,g^x)$ which represents its public key.
    \item $A$ sends $A'$ two messages $m_0,m_1$ and $A'$ flips a bit $b$ and returns ciphertext $(g^y, g^c \cdot m_b)$.
    \item $A'$ outputs $1$ if $A$ outputs $b$.
\end{enumerate}
Here, if it is the case that $c = xy$ then $A$'s view is exactly the El gamal protocol, so $A'$ guesses that $g^c = g^{xy}$ correctly with probability $\epsilon$. In the other case, $g^c \cdot m_b$ is a completely random element in the group, so the probability that it gets the bit correct in this case is exactly $1/2$. Since the probability that the algorithm outputs $1$ in each case differ by a non-negligible amount this implies that $A'$ can differentiate between the two worlds with non-negligible advantage giving us a contradition.

\end{document}